\textbf{Problem 5.13 (Dujella, 2026).} For a positive integer $k$, denote by $n_1(k)$ the smallest (in the absolute value) non-zero integer for which there exists a triple $\{a,b,c\}$ of non-zero integers and a set $N$ of integers, such that $|N|=k$, $n_1(k) \in N$, and $\{a,b,c\}$ is a $D(n)$-triple for all $n \in N$.
\begin{enumerate}
\item Is it true that $n_1(5)=36$?
\item What can be said (assuming some plausible conjectures) about the asymptotic behaviour of $|n_1(k)|$ for large $k$?
\end{enumerate}

\medskip

\textit{Remark.} Sadowski (2026) answered question (a) showing that $n_1(5) \leqslant 4$. In fact, he showed that also $n_1(6) \leqslant 4$ with the following example: $\{a,b,c\}=\{30,1056,65520\}$, $N=\{4,266436,15248601,21479364,63936324,1037903409\}$. Furthermore, $n_1(7) \leqslant 36$ with $\{a,b,c\}=\{8,1620,257040\}$, $N=\{36,857529,1470564,200322756,238596849,452017161,16308809476\}$.
